% \documentclass{ximera}
% \typeout{Start loading xmPreamble.tex}%

% Add here extra macro's that are loaded automatically by all documents of claas 'ximera' or 'xourse' in this repo

%%
%%  Example:
%%
% \newcommand{\R}{\mathbb{R}}
\usepackage{tikz}
\usetikzlibrary{positioning, arrows.meta}
\usepackage{multicol}
\usepackage{enumitem}
\usepackage{caption}
\usepackage{amsmath}
\usepackage{amsthm}


% \title{Markov Chains and Processes On Discrete State Spaces}
% \begin{document}
% \begin{abstract}
% \end{abstract}
% \maketitle
% \begin{theorem}[5.2 Rules for Exponents]
% For any numbers $b$, $c$, $x$, and $y$ for which the expressions that follow are defined, 
% \begin{tasks}[label=(\alph*), label-width=4ex](3)
% \task $\displaystyle b^{x+y} = b^xb^y$
% \task $\displaystyle (bc)^x = b^xc^x$
% \task $\displaystyle b^{xy} = (b^x)^y = (b^y)^x$
% \task $\displaystyle b^{x-y} = \dfrac{b^x}{b^y}$
% \task $\left( \dfrac{b}{c} \right)^x = \dfrac{b^x}{c^x}$
% \task $\displaystyle b^{x/y} = (b^x)^{1/y} = (b^{1/y})^x$
% \end{tasks}
% \end{theorem}
% \end{document}

\documentclass{ximera}
\typeout{Start loading xmPreamble.tex}%

% Add here extra macro's that are loaded automatically by all documents of claas 'ximera' or 'xourse' in this repo

%%
%%  Example:
%%
% \newcommand{\R}{\mathbb{R}}
\usepackage{tikz}
\usetikzlibrary{positioning, arrows.meta}
\usepackage{multicol}
\usepackage{enumitem}
\usepackage{caption}
\usepackage{amsmath}
\usepackage{amsthm}

\title{Rules for Exponents}
\begin{document}
\begin{abstract}
\end{abstract}
\maketitle
\begin{theorem}[5.2 Rules for Exponents]
For any numbers $b$, $c$, $x$, and $y$ for which the expressions that follow are defined,
\begin{multicols}{3}
\begin{enumerate}[label=(\alph*)]
\item $\displaystyle b^{x+y} = b^xb^y$
\item $\displaystyle (bc)^x = b^xc^x$
\item $\displaystyle b^{xy} = (b^x)^y = (b^y)^x$
\item $\displaystyle b^{x-y} = \dfrac{b^x}{b^y}$
\item $\left( \dfrac{b}{c} \right)^x = \dfrac{b^x}{c^x}$
\item $\displaystyle b^{x/y} = (b^x)^{1/y} = (b^{1/y})^x$
\end{enumerate}
\end{multicols}
\end{theorem}
\end{document}